\documentclass[../main.tex]{subfiles}

\begin{document}
% -----------------------------
\section{Front end}
% -----------------------------
So far, we have been constructing modules manually using our HC dialect operations.
In this step, we will move toward a more realistic workflow by loading a model from
an ONNX file instead of building the IR by hand. Since we do not yet have
an ONNX model prepared, we will first create a small example model and export
it to ONNX format. Then we will implement a simple front end that reads the ONNX
graph and maps its nodes to our corresponding HC operations, producing an HC module
that can be further compiled.

To organize this part of the project, we will create a folder with the following structure:
\begin{lstlisting} [language=bash, caption={Front end folder structure}]
front_end
├── __init__.py
├── build_model.py
├── loader.py
\end{lstlisting}

In \path{build_model.py}, import \texttt{onnx} related packages and define final model path:
\begin{lstlisting} [language=Python, caption={onnx imports}]
import onnx
from onnx import helper, TensorProto

MODEL_PATH = "score_model.onnx"
\end{lstlisting}


\begin{lstlisting} [language=Python, caption={Building model}]
def build_score_model(path: str = MODEL_PATH):
    # ------------------------------------------------------------
    # Parameters
    # ------------------------------------------------------------
    MU = 40.0
    ALPHA = 2.0
    P = 2.0
    LO = 0.0
    HI = 1000.0

    # ------------------------------------------------------------
    # Input / Output
    # ------------------------------------------------------------
    x = helper.make_tensor_value_info("x", TensorProto.FLOAT, [])
    y = helper.make_tensor_value_info("score", TensorProto.FLOAT, [])

    # ------------------------------------------------------------
    # Constant initializers
    # ------------------------------------------------------------
    mu_init = helper.make_tensor("mu", TensorProto.FLOAT, [], [MU])
    a_init  = helper.make_tensor("alpha", TensorProto.FLOAT, [], [ALPHA])
    p_init  = helper.make_tensor("p", TensorProto.FLOAT, [], [P])
    lo_init = helper.make_tensor("lo", TensorProto.FLOAT, [], [LO])
    hi_init = helper.make_tensor("hi", TensorProto.FLOAT, [], [HI])

    # ------------------------------------------------------------
    # Graph nodes
    # ------------------------------------------------------------
    nodes = [
        # d1 = x - mu
        helper.make_node("Sub", ["x", "mu"], ["d1"], name="sub_x_mu"),

        # d2 = mu - x
        helper.make_node("Sub", ["mu", "x"], ["d2"], name="sub_mu_x"),

        # abs_d = max(d1, d2)
        helper.make_node("Max", ["d1", "d2"], ["abs_d"], name="abs_max"),

        # e = pow(abs_d, p)
        helper.make_node("Pow", ["abs_d", "p"], ["e"], name="pow"),

        # s = alpha * e
        helper.make_node("Mul", ["alpha", "e"], ["s"], name="scale"),

        # lower = max(s, lo)
        helper.make_node("Max", ["s", "lo"], ["lower"], name="clamp_low"),

        # clamped = min(lower, hi)
        helper.make_node("Min", ["lower", "hi"], ["clamped"], name="clamp_high"),

        # score = relu(clamped)
        helper.make_node("Relu", ["clamped"], ["score"], name="relu"),
    ]

    graph = helper.make_graph(
        nodes=nodes,
        name="ScoreGraph",
        inputs=[x],
        outputs=[y],
        initializer=[mu_init, a_init, p_init, lo_init, hi_init],
    )

    model = helper.make_model(
        graph,
        opset_imports=[helper.make_opsetid("", 13)],
        producer_name="score_builder",
    )

    onnx.checker.check_model(model)
    onnx.save(model, path)
    print(f"Saved ONNX model to: {path}")


if __name__ == "__main__":
    build_score_model()
\end{lstlisting}
Run the script and observe generated model using \path{netron.app} website.
Then, in \path{loader.py}, we will build HC module based on provided onnx model.
Start with importing required packages and our dialect:
\begin{lstlisting} [language=Python, caption={Loader imports}]
import onnx
from xdsl.dialects import func, arith

from xdsl.context import Context
from xdsl.dialects.builtin import ModuleOp, i32
from xdsl.ir import Region, Block

from hc_dialect import HiCompiler, HCAdd, HCSub, HCMul, HCRelu, HCPow, HCMax, HCMin
\end{lstlisting}
Define helper functions for types conversion:
\begin{lstlisting} [language=Python, caption={Defining loader helper functions}]
# --------------------------------------
#  Helper functions
# --------------------------------------
def _onnx_tensor_to_scalar_int(tensor_proto) -> int:
    arr = onnx.numpy_helper.to_array(tensor_proto)
    # Restriction: scalar or 1-element tensor
    if arr.size != 1:
        raise ValueError(f"Only scalar/1-element constants supported, got shape {arr.shape}")
    return int(arr.reshape(()))

def _make_i32_const(value: int) -> arith.ConstantOp:
    return arith.ConstantOp.from_int_and_width(int(value), 32)
\end{lstlisting}

\begin{lstlisting} [language=Python, caption={Defining model loader}]
# --------------------------------------
#  Loader
# --------------------------------------
def import_onnx_to_hc_module(
    ctx: Context, onnx_path: str, fn_name: str="main"
) -> ModuleOp:
    model = onnx.load(onnx_path)
    graph = model.graph

    # Map ONNX value names -> SSAValue (results)
    env: dict[str, object] = {}

    # Assume one scalar input
    # Map ONNX graph input name to function/block argument
    entry_block = Block(arg_types=[i32])
    x_name = graph.input[0].name
    env[x_name] = entry_block.args[0]

    ops = []

    # 1) Initializers (weights/constants stored in graph.initializer)
    for init in graph.initializer:
        v = _onnx_tensor_to_scalar_int(init)
        c = _make_i32_const(v)
        ops.append(c)
        env[init.name] = c.result

    # 2) Constant nodes (op_type == "Constant")
    # ONNX Constant usually stores a tensor in attribute "value"
    def emit_constant_node(node):
        value_attr = None
        for a in node.attribute:
            if a.name == "value":
                value_attr = a.t
        if value_attr is None:
            raise ValueError("Constant node without 'value' attribute")
        v = _onnx_tensor_to_scalar_int(value_attr)
        c = _make_i32_const(v)
        ops.append(c)
        env[node.output[0]] = c.result

    # 3) Node lowering
    for node in graph.node:
        if node.op_type == "Constant":
            emit_constant_node(node)
            continue

        def get(name: str):
            if name not in env:
                raise KeyError(f"ONNX value not found yet: {name} (node {node.op_type})")
            return env[name]

        if node.op_type == "Add":
            a = get(node.input[0])
            b = get(node.input[1])
            hc_add = HCAdd(operands=[a, b], result_types=[i32])
            ops.append(hc_add)
            env[node.output[0]] = hc_add.results[0]
            continue

        if node.op_type == "Sub":
            a = get(node.input[0])
            b = get(node.input[1])
            hc_sub = HCSub(operands=[a, b], result_types=[i32])
            ops.append(hc_sub)
            env[node.output[0]] = hc_sub.results[0]
            continue

        if node.op_type == "Mul":
            a = get(node.input[0])
            b = get(node.input[1])
            hc_mul = HCMul(operands=[a, b], result_types=[i32])
            ops.append(hc_mul)
            env[node.output[0]] = hc_mul.results[0]
            continue

        if node.op_type == "Relu":
            x = get(node.input[0])
            hc_relu = HCRelu(operands=[x], result_types=[i32])
            ops.append(hc_relu)
            env[node.output[0]] = hc_relu.results[0]
            continue

        if node.op_type == "Pow":
            a = get(node.input[0])
            b = get(node.input[1])
            hc_pow = HCPow(operands=[a, b], result_types=[i32])
            ops.append(hc_pow)
            env[node.output[0]] = hc_pow.results[0]
            continue

        if node.op_type == "Max":
            a = get(node.input[0])
            b = get(node.input[1])
            hc_max = HCMax(operands=[a, b], result_types=[i32])
            ops.append(hc_max)
            env[node.output[0]] = hc_max.results[0]
            continue

        if node.op_type == "Min":
            a = get(node.input[0])
            b = get(node.input[1])
            hc_min = HCMin(operands=[a, b], result_types=[i32])
            ops.append(hc_min)
            env[node.output[0]] = hc_min.results[0]
            continue

        raise NotImplementedError(f"Unsupported ONNX op: {node.op_type}")

    # 4) Determine outputs (graph.output)
    if len(graph.output) != 1:
        raise ValueError("For now: expect exactly 1 model output")
    out_name = graph.output[0].name
    if out_name not in env:
        raise KeyError(f"Graph output {out_name} not produced")
    ret = func.ReturnOp(env[out_name])
    ops.append(ret)

    for op in ops:
        entry_block.add_op(op)

    fn = func.FuncOp(fn_name, ([i32], [i32]))
    fn.body = Region(entry_block)

    return ModuleOp(ops=[fn])
\end{lstlisting}
This function converts a neural-network model stored in the ONNX format into an
intermediate representation (IR) expressed using the HC dialect in xDSL. The goal
is to transform the graph-based ONNX model into a sequence of SSA operations that
can be further processed, optimized, or executed by your compiler infrastructure.

The loader reads the ONNX graph and builds an internal mapping between ONNX value
names and SSA values in the IR. Model inputs are represented as function arguments,
while constants and parameters are turned into constant operations. Each ONNX node
is then translated into an equivalent operation in the HC dialect, gradually constructing
the computation inside the function’s entry block.

Once all nodes are translated, the loader determines the graph’s output and adds a
return operation to produce the final value. The resulting operations are placed into
a function, and that function is wrapped inside a module, producing a complete IR
representation of the original ONNX model that can be executed or transformed later
in the compilation pipeline.

Since the module we just build has one argument that we can pass as an input, we need to
modify our \path{interpreter.py} to also support this feature by changing \path{run\_module}
signature and providing logic for args handling:
\begin{lstlisting} [language=Python, caption={Argument handling in interpreter}]
def run_module(self, module,  args: list[int], func_name: str = "my_func") -> int:
...
                    if name == func_name or (name is None and func_name == "my_func"):
                        # assume single-block body for now
                        body_block = list(op.body.blocks)[0]

                        if len(args) != len(body_block.args):
                            raise RuntimeError(
                                f"Expected {len(body_block.args)} args, got {len(args)}"
                            )
                        for block_arg, value in zip(body_block.args, args):
                            self._set(block_arg, value)
                        return self.run_block(body_block)
\end{lstlisting}

Now, in our \path{hc_main.py}, we can entirely remove \path{build_module} function and use
loader.

\begin{lstlisting} [language=Python, caption={hc\_main.py imports}]
from front_end.build_model import MODEL_PATH
from front_end.loader import import_onnx_to_hc_module
\end{lstlisting}
Load the model and run it through interpreter.
\begin{lstlisting} [language=Python, caption={Using loader in main}]
def main():
    ctx = build_context()
    module = import_onnx_to_hc_module(ctx, MODEL_PATH, "my_func")

    print("=== HIGH LEVEL (hc.*) ===")
    print(module)
    input_val = [31]

    hc_interpreter = Interpreter()
    result_before = hc_interpreter.run_module(module, args=input_val, func_name="my_func")
    print("\nResult (interpreted, before lowering):", result_before)

    config = MiddleEndPipelineConfig(
        apply_lowering=True,
        apply_constant_folding=False,
        apply_dce=False,
        run_analysis=False,
        debug_mode=False,
    )

    middle_end_pipeline = MiddleEndPipeline(config)
    middle_end_pipeline.apply_passes(module)
    print("=== After lowering ===")
    print(module)

    lower_interpreter = Interpreter()
    result_after = lower_interpreter.run_module(module, args=input_val, func_name="my_func")
    print("\nResult (interpreted, after lowering):", result_after)
\end{lstlisting}
Try different values of \path{input_val}, and observe how output changes when we diverge from
the value 40.

\end{document}
